\section{What is the optimal policy?}
\label{sec:optimal-control}
For known parameters, minimizing the full-state regret coefficient is equivalent to maximizing causal average reward. Since the states are stationary and exogenous,
\begin{equation}
c_\pi=G-\liminf_{T\to\infty}\frac1T\E_\pi\sum_{t=1}^TX_{A_t,t}.
\label{eq:gain-equivalence}
\end{equation}
The current beliefs $b=((m_i,v_i))_{i=1}^K$ are a sufficient state for this control problem. Observing $x$ from arm $i$ produces the next-round belief $\mathcal T_i(b,x)$ given by
\begin{align}
m_i'&=\mu_i+\phi_i(x-\mu_i), &v_i'&=q_i,\nonumber\\
m_j'&=\mu_j+\phi_j(m_j-\mu_j), &v_j'&=\phi_j^2v_j+q_j\quad(j\ne i).
\label{eq:belief-transition}
\end{align}
Thus only the selected arm's next predictive mean is random at this transition. Unselected arms still accumulate uncertainty.

\begin{proposition}[Exact finite-horizon control]
Let $J_n(b)$ be the largest expected total reward over $n$ remaining rounds from belief $b$. Then $J_0=0$ and
\begin{equation}
J_n(b)=\max_i\left\{m_i+\E_{x\sim N(m_i,v_i)}J_{n-1}(\mathcal T_i(b,x))\right\}.
\label{eq:finite-bellman}
\end{equation}
Selecting a maximizing arm at each remaining horizon is optimal for that finite-horizon problem.
\end{proposition}
\begin{proof}
Given the belief and first action $i$, the current expected reward is $m_i$ and the observed reward has distribution $N(m_i,v_i)$. Equation~\eqref{eq:belief-transition} is its complete next-round sufficient statistic. Conditional optimization of the remaining $n-1$ actions followed by maximization over the first action proves the recursion by induction.
\end{proof}
This is an exact policy characterization, not a computed solution for the experimental cases. Each action expectation is one-dimensional, but its continuation value depends on the joint belief of every arm. Recursive integration and representation of that joint value are the main computational obstacles.

The information bonus can also be isolated explicitly. Let $b^-$ be the hypothetical next-round belief if no current state were observed: every mean and variance undergo the unselected-arm update in~\eqref{eq:belief-transition}. Then
\[
I_{i,n}(b)=\E J_{n-1}(\mathcal T_i(b,X_i))-J_{n-1}(b^-),
\qquad A\in\argmax_i\{m_i+I_{i,n}(b)\}.
\]
Here $I_{i,n}\geq0$: a continuation policy can always ignore the new observation and reproduce the no-observation continuation. Unlike a variance bonus, this information value depends on decision gaps and the opportunity to obtain further observations. For $n=1$ it is zero; for $n=2$ it recovers the exact two-period information value.

For the linear coefficient, the corresponding average-reward optimality equation is
\begin{equation}
g+h(b)=\max_i\left\{m_i+\E h(\mathcal T_i(b,X_i))\right\},
\qquad X_i\sim N(m_i,v_i).
\label{eq:average-bellman}
\end{equation}
Here $h$ measures the future advantage of a belief relative to the average reward $g$. A maximizing selector is average optimal if a solution has the transversality property
\begin{equation}
\frac{\E_\pi|h(B_{T+1})|}{T}\longrightarrow0
\label{eq:transversality}
\end{equation}
for all admissible policies from our initial belief. The next section proves this assertion by telescoping. We do not establish existence of such a solution in the unbounded Gaussian belief space. Writing the equation alone is therefore not an average-optimality theorem for an implemented algorithm.

The continuation term is the appropriate information bonus: it values an observation in the context of competitors, their uncertainties, and subsequent observations. Positive persistence alone does not justify a universally optimal independent-arm index. General restless Markov counterexamples already invalidate universal optimality of independent-arm indices~\cite{ortner2014}; those examples do not prove nonexistence of an optimal index for this particular Gaussian family. A Whittle policy would require a separate indexability argument and performance analysis.

\section{Certifying the regret coefficient of a candidate}
\label{sec:certificate}
An approximate continuation function can be useful even without solving~\eqref{eq:average-bellman}. For a trial function $h$, define
\begin{equation}
D_h(b,i)=m_i+\E h(\mathcal T_i(b,X_i))-h(b),\qquad
r_h(b)=\max_iD_h(b,i),
\end{equation}
and let $\pi_h$ choose a maximizing arm. The residual $r_h$ is an average-reward Bellman residual, rather than an empirical regret estimate.

\begin{proposition}[Residual certificate]
Suppose $h$ satisfies~\eqref{eq:transversality} for all admissible policies, all expectations above are finite, and
$\ell\leq r_h(b)\leq u$ on every reachable belief. Then
\begin{equation}
G-u\leq c^*\leq c_{\pi_h}\leq G-\ell,\qquad
c_{\pi_h}-c^*\leq u-\ell.
\label{eq:certificate}
\end{equation}
In particular, an exact constant residual certifies optimality.
\end{proposition}
\begin{proof}
For any policy, the conditional expected reward plus expected change in $h$ is at most $r_h(B_t)\leq u$. Summing gives
\[
\E_\pi\sum_{t=1}^TX_{A_t,t}\leq Tu+h(B_1)-\E_\pi h(B_{T+1}).
\]
For $\pi_h$ the conditional expression equals $r_h(B_t)\geq\ell$, yielding the reverse inequality with $T\ell$. Divide by $T$, use transversality, and apply~\eqref{eq:gain-equivalence}. Infimizing over policies and subtracting the two bounds proves~\eqref{eq:certificate}.
\end{proof}

The Gaussian setting also admits a useful weighted version. Write $W(b)=\sum_i|m_i-\mu_i|$. Under stationary exogenous states, for every policy and round,
\begin{equation}
\E_\pi W(B_t)\leq\sqrt{2/\pi}\sum_i\sqrt{V_i}=:C_V,
\label{eq:belief-moment}
\end{equation}
by $m_i=\E[X_{i,t}\mid H_t]$ and conditional Jensen's inequality. Consequently $|h(b)|\leq C_0+C_1W(b)$ suffices for transversality. If a globally verified residual bound has the form
\begin{equation}
|r_h(b)-g|\leq\epsilon_0+\epsilon_1W(b),
\end{equation}
the same telescoping argument gives, with $\epsilon=\epsilon_0+\epsilon_1C_V$,
\begin{equation}
G-g-\epsilon\leq c^*\leq c_{\pi_h}\leq G-g+\epsilon,
\qquad c_{\pi_h}-c^*\leq2\epsilon.
\end{equation}
This allows unbounded means while bounding average error. A residual checked only on sampled beliefs or a finite grid is not a global certificate; interpolation, integration errors, and omitted Gaussian tails also need control. No numerical certificate is claimed here.

There is a related conditional improvement guarantee. If a baseline $\pi_0$ has gain $g_0$ and a bias $h_0$ satisfying its policy evaluation equation
\[
g_0=\sum_i\pi_0(i\mid b)D_{h_0}(b,i),
\]
then $r_{h_0}(b)\geq g_0$. Under the same transversality conditions, greedifying with respect to $h_0$ gives $c_{\pi_{h_0}}\leq G-g_0=c_{\pi_0}$. This suggests evaluating PS's bias and improving its actions, with a separate residual certificate for closeness to the optimum. A fitted bias or a truncated Monte Carlo rollout does not inherit this guarantee without error control. A discounted-return improvement alone also does not certify improvement of our average-reward objective.

\section{A lifetime-weighted information rule and its limit}
\label{sec:long-kg}
To test whether a simple long-horizon correction improves the two-period policy, consider common mean $\mu$ and common $\phi\in[0,1)$. Put $C_i=\max_{j\ne i}m_j$, $d_i=|m_i-C_i|$, and define the expected improvement in the largest predictive mean after one exact observation by
\begin{align}
\mathrm{KG}_i(b)
&=\E\max\{C_i,X_i\}-\max_jm_j\nonumber\\
&=\sqrt{v_i}\varphi\left(\frac{d_i}{\sqrt{v_i}}\right)
-d_i\Phi\left(-\frac{d_i}{\sqrt{v_i}}\right).
\label{eq:kg}
\end{align}
Knowledge-gradient measurement policies are established in Bayesian information collection~\cite{frazier2008}; their ranking-and-selection guarantees do not apply to this restless reward problem.

Suppose the continuation retains this one observation but ignores every later observation. At lag $s$, all predictive means share the same affine contraction toward $\mu$. The expected improvement in the best predictive mean is therefore exactly $\phi^s\mathrm{KG}_i(b)$. The cumulative improvement is
\begin{equation}
\sum_{s=1}^\infty\phi^s\mathrm{KG}_i(b)
=\frac\phi{1-\phi}\mathrm{KG}_i(b).
\end{equation}
This continuation calculation is finite for $\phi<1$ and yields the illustrative action rule
\begin{equation}
A_t\in\argmax_i\left\{m_i+\frac\phi{1-\phi}\mathrm{KG}_i(b)\right\}.
\label{eq:long-kg}
\end{equation}
It is deterministic, accounts for decision gaps as well as variance, and has no fitted multiplier. The information term vanishes at $\phi=0$. For common parameters the two-period rule uses $\phi\mathrm{KG}_i$ instead; at $\phi=0.99$ the lifetime multiplier is $99$, compared with $0.99$ for the two-period rule. Predictive mean improvements decay as $\phi^s$, whereas prediction-variance improvements decay as $\phi^{2s}$: their lifetime factors measure different quantities.

Equation~\eqref{eq:long-kg} is not a solution of~\eqref{eq:average-bellman}. Replanning after every observation uses information that the continuation calculation discarded, so it can count benefits that subsequent observations would replace. Heterogeneous means or persistence also break its common-contraction simplification.

The additional comparison reuses the original six homogeneous cases, all 40 independent trajectories per case, and the same burn-in and measured horizon. An exact replay of a saved PS replication verifies that running the new policy separately preserves the common state trajectory. No multiplier is selected from the results.
\begin{table}[ht]
\centering
\caption{Additional common-trajectory comparison. Smaller coefficients are better; positive paired PS minus KG differences favor KG. Intervals are approximate 95\% intervals.}
\resizebox{\textwidth}{!}{\begin{tabular}{rrrrrr}
\toprule
$\phi$ & Greedy & Two-step & PS & Long-horizon KG & PS minus KG\\
\midrule
\inputtable{results/policy_improvement/results_table.tex}
\bottomrule
\end{tabular}}
\end{table}
At $\phi=0.99$, long-horizon KG has measured coefficient $0.1659\pm0.0011$, higher than PS's $0.1495\pm0.0010$. This is a paired comparison on the same exogenous trajectories. It does not establish optimality of either policy.

The lifetime rule improves on PS at moderate persistence, but two-period control remains stronger at $\phi=0.9$ and PS remains stronger at $\phi=0.99$ and $0.995$. These results rule out presenting this geometric bonus as a uniformly better policy. They motivate estimating the actual average-reward continuation value, whose information benefit accounts for future refreshes and competing observations.
